\documentclass[12pt]{article}
\usepackage[margin=0.3in,top=0.25in,bottom=0.4in]{geometry}
\makeatletter
\renewcommand{\@maketitle}{\newpage\null\vskip -3.5em\begin{center}{\LARGE \@title \par}\end{center}\par\vskip 0.2em}
\makeatother
\usepackage{amsmath,amssymb}
\usepackage{booktabs}
\usepackage{hyperref}
\setlength{\parskip}{0.6em}
\setlength{\parindent}{0pt}

\title{The BCFY fragmentation function was too large}
\author{}
\date{}

\begin{document}
\maketitle

\section*{In short}

Our BCFY fragmentation function was \textbf{2.4 times too large}, because it was not normalized.
All our BCFY results were therefore too large by about this factor. This is why BCFY was always
2--3 times above Kniehl--Kramer and HymnD. Kniehl--Kramer and HymnD are correct.

\section*{1. What a fragmentation function must satisfy}

A fragmentation function $D(z)$ tells how a charm quark turns into a $D^0$ meson that carries a
fraction $z$ of its momentum. Its integral,
\[
  \int_0^1 dz\,D(z),
\]
is the \textbf{probability} that the charm quark ends up as a $D^0$. A probability cannot be larger
than 1, and the measured value is about $0.56$.

Our functions give:
\begin{center}
\begin{tabular}{lc}
\toprule
 & $\int dz\,D(z)$ \\
\midrule
Kniehl--Kramer & 0.74 \\
HymnD          & 0.64 \\
BCFY (before)  & \textbf{1.36} \\
\bottomrule
\end{tabular}
\end{center}
A probability of 1.36 is impossible, so something is wrong with BCFY.

\section*{2. How BCFY is built}

A charm quark can become a $D^0$ in two ways: directly, or first as a $D^*$ that then decays into a
$D^0$. Following Cacciari \& Nason (JHEP 09 (2003) 006), the BCFY function is
\[
  D(z) = 0.168 \times \big[\text{direct shape}\big] \;+\; 0.39 \times \big[\text{shape from } D^*\big].
\]
The numbers $0.168$ and $0.39$ are the measured probabilities of the two ways, and they add up to
$0.558$. The two shapes, taken from the BCFY paper, only describe \emph{how} the momentum is shared.
For the formula to make sense, each shape must have integral 1. Cacciari \& Nason say this
explicitly: the shapes are ``understood to be normalized to one''.

\section*{3. What went wrong}

We used the two shapes exactly as written in the BCFY paper, where they are not normalized. For
$r=0.1$ their integrals are $1.507$ and $2.825$ instead of 1:
\begin{center}
\begin{tabular}{lccc}
\toprule
 & direct & from $D^*$ & total \\
\midrule
before (wrong)  & $0.168\times1.507=0.25$ & $0.39\times2.825=1.10$ & $1.36$ \\
now (correct)   & $0.168$ & $0.39$ & $0.558$ \\
\bottomrule
\end{tabular}
\end{center}
The total was too large by $1.36/0.558\approx2.4$.

The same unnormalized formula is in:
\begin{itemize}
\item \texttt{EKO-FF} (\texttt{d0-fragmentation/parametrizations.py}), which made the BCFY grid used
  in this project;
\item the older codes \texttt{DiffractiveD0/inclusive}, \texttt{D0CS}, \texttt{flux2d} and
  \texttt{inclusive\_test}.
\end{itemize}

\section*{4. The fix}

Each shape is divided by its own integral, so that it has integral 1, before multiplying by $0.168$
and $0.39$. Then the BCFY function integrates to $0.558$, as it should. The BCFY grid has been
regenerated with this fix (same leading-order evolution as before).

\section*{5. What this changes}

\begin{itemize}
\item All BCFY cross sections go down by a factor of about 2.4 (the exact factor depends slightly on
  the kinematics, because the two parts are rescaled by different amounts).
\item BCFY will then be much closer to Kniehl--Kramer and HymnD. It remains somewhat larger, because
  its shape is harder (it peaks at larger $z$).
\item Kniehl--Kramer and HymnD do not change.
\item All BCFY results need to be recomputed.
\end{itemize}

\section*{6. Check of the new grid}

The new grid (\texttt{input/BCFY\_EKO/bcfy\_eko\_0000.dat}; the old one is kept as
\texttt{bcfy\_eko\_0000\_unnormalized.dat}) uses the same $z_h$ and $Q$ points as before
($200\times30$, $Q=1.5$--$50$~GeV). Its integral is $0.559$--$0.561$ at all scales, as it should
be (the evolution is the same leading-order DGLAP as before).

To see the effect on a cross section, the table compares $\int dz\,z^4D(z)$, which is roughly
what a spectrum falling like $p_\perp^{-5}$ picks up:
\begin{center}
\begin{tabular}{cccc}
\toprule
$Q$ [GeV] & new / old BCFY & new BCFY / Kniehl--Kramer & new BCFY / HymnD \\
\midrule
1.5  & 0.42 & 1.41 & 1.44 \\
3.35 & 0.42 & 1.41 & 1.53 \\
9.12 & 0.42 & 1.41 & 1.58 \\
50   & 0.42 & 1.41 & 1.61 \\
\bottomrule
\end{tabular}
\end{center}
So the BCFY results go down by a factor of about 2.4, and BCFY ends up about 1.4--1.6 times
Kniehl--Kramer and HymnD, instead of 2--3 times. The remaining difference comes from its harder
shape.

\end{document}
